\section{Existence of cluster points}

Throughout, $S$ is the parameter set, $(E, \mathcal{E})$ the state space, $\Cfg = E^S$,
$\mathcal{F} = \mathcal{E}^S$, $\mathcal{F}_\Lambda$ the $\sigma$-algebra of events depending
only on the coordinates in $\Lambda$,
$\mathcal{F}^0 = \bigcup_{\Lambda \in \Fin} \mathcal{F}_\Lambda$ the algebra of local events, and
$\Prob(\Cfg, \mathcal{F})$ carries the topology of local convergence of \S4.1. Families
$(\mu_a)_{a \in \iota}$ are indexed along a filter $l$ on $\iota$; a net is the case of the
filter of a directed set.

\subsection*{Equicontinuity}

\begin{definition}[Local equicontinuity, Georgii (4.6)]
    \label{def:cp-loc-equi}
    \uses{def:cylinder-event}
    \lean{MeasureTheory.GibbsMeasure.LocallyEquicontinuous}
    \leanok

    A family $(\mu_a)_{a \in \iota}$ in $\Prob(\Cfg, \mathcal{F})$, indexed along a filter $l$,
    is \textbf{locally equicontinuous} if for each $\Lambda \in \Fin$ and each sequence
    $(A_m)_{m \geq 1}$ in $\mathcal{F}_\Lambda$ with $A_m \downarrow \emptyset$,
    \[\lim_{m \to \infty} \limsup_{a \in l} \mu_a(A_m) = 0.\]
\end{definition}

\begin{definition}[Equicontinuity, Georgii (4.5)]
    \label{def:cp-equi}
    \lean{MeasureTheory.GibbsMeasure.EquicontinuousOnLocalEvents}
    \leanok

    The family $(\mu_a)_{a \in \iota}$ is \textbf{equicontinuous} if
    $\lim_{m \to \infty} \limsup_{a \in l} \mu_a(A_m) = 0$ for every sequence
    $(A_m)_{m \geq 1}$ in $\mathcal{F}^0$ with $A_m \downarrow \emptyset$.
\end{definition}

\begin{lemma}[Georgii (4.5) implies (4.6)]
    \label{lem:cp-equi-loc-equi}
    \uses{def:cp-loc-equi, def:cp-equi}
    \lean{MeasureTheory.GibbsMeasure.EquicontinuousOnLocalEvents.locallyEquicontinuous}
    \leanok

    An equicontinuous family is locally equicontinuous.
\end{lemma}
\begin{proof}
    \leanok

    $\mathcal{F}_\Lambda \subseteq \mathcal{F}^0$.
\end{proof}

\begin{theorem}[Necessary condition for a cluster point, Georgii (4.4)]
    \label{thm:cp-necessary}
    \lean{MeasureTheory.GibbsMeasure.tendsto_liminf_zero_of_mapClusterPt}
    \leanok

    If $\mu$ is a cluster point of $(\mu_a)_{a \in \iota}$ along $l$ in the topology of local
    convergence, then for every sequence $(A_m)_{m \geq 1}$ in $\mathcal{F}^0$ with
    $A_m \downarrow \emptyset$,
    \[\lim_{m \to \infty} \liminf_{a \in l} \mu_a(A_m) = 0.\]
\end{theorem}
\begin{proof}
    \leanok

    Take an ultrafilter $U \leq l$ along which $\mu_a \to \mu$. For each $m$,
    $\liminf_{a \in l} \mu_a(A_m) \leq \liminf_{a \in U} \mu_a(A_m) = \mu(A_m)$, and
    $\mu(A_m) \to 0$ by $\sigma$-additivity.
\end{proof}

\begin{definition}[Standard Borel space, Georgii (4.7)]
    \label{def:cp-standard-borel}
    \lean{StandardBorelSpace}
    \leanok

    A measurable space $(E, \mathcal{E})$ is a \textbf{standard Borel space} if $\mathcal{E}$
    is the Borel $\sigma$-algebra of some complete separable metric on $E$; equivalently, of
    some Polish topology on $E$.
\end{definition}

\subsection*{Proof of Georgii (4.9): ultrafilter limit, marginals, Kolmogorov}

\begin{definition}[Georgii's finitely additive limit $\mu^0$]
    \label{def:cp-ultrafilter-limit}
    \lean{MeasureTheory.GibbsMeasure.ultrafilterLimit,
      MeasureTheory.GibbsMeasure.tendsto_ultrafilterLimit}
    \leanok

    For an ultrafilter $U$ on $\iota$ and $A \subseteq \Cfg$ put
    $\mu^0(A) = \lim_{a \in U} \mu_a(A) \in [0,\infty]$; the limit exists by compactness of
    $[0,\infty]$.
\end{definition}

\begin{lemma}[$\mu^0$ is $\sigma$-additive on each $\mathcal{F}_\Lambda$]
    \label{lem:cp-sigma-additive}
    \uses{def:cp-loc-equi, def:cp-ultrafilter-limit}
    \lean{MeasureTheory.GibbsMeasure.ultrafilterLimit_union,
      MeasureTheory.GibbsMeasure.ultrafilterLimit_le_limsup,
      MeasureTheory.GibbsMeasure.tendsto_ultrafilterLimit_zero,
      MeasureTheory.GibbsMeasure.ultrafilterLimit_iUnion}
    \leanok

    Let $U \leq l$ be an ultrafilter and $(\mu_a)_{a \in \iota}$ locally equicontinuous along
    $l$. Then $\mu^0$ is finitely additive, $\mu^0(A) \leq \limsup_{a \in l} \mu_a(A)$ for all
    $A$, and for each $\Lambda \in \Fin$, $\mu^0$ is $\sigma$-additive on $\mathcal{F}_\Lambda$.
\end{lemma}
\begin{proof}
    \leanok

    Finite additivity is continuity of addition on $[0,\infty]$; the bound holds because
    $U \leq l$. If $A_m \in \mathcal{F}_\Lambda$ and $A_m \downarrow \emptyset$ then
    $\lim_m \mu^0(A_m) \leq \lim_m \limsup_{a \in l} \mu_a(A_m) = 0$ by local equicontinuity,
    and a finitely additive content on a set ring that is continuous at $\emptyset$ is
    $\sigma$-additive.
\end{proof}

\begin{definition}[Finite-volume marginals]
    \label{def:cp-finvol}
    \uses{lem:cp-sigma-additive}
    \lean{MeasureTheory.GibbsMeasure.finVolMeasure,
      MeasureTheory.GibbsMeasure.finVolMeasure_apply,
      MeasureTheory.GibbsMeasure.isProbabilityMeasure_finVolMeasure}
    \leanok

    For $\Lambda \in \Fin$, the map $B \mapsto \mu^0(\{\sigma_\Lambda \in B\})$,
    $B \in \mathcal{E}^\Lambda$, is a probability measure on $E^\Lambda$, the marginal of
    $\mu^0$ in $\Lambda$.
\end{definition}

\begin{lemma}[Consistency and Kolmogorov extension]
    \label{lem:cp-projective}
    \uses{def:cp-finvol, def:cp-standard-borel}
    \lean{MeasureTheory.GibbsMeasure.isProjectiveMeasureFamily_finVolMeasure,
      MeasureTheory.GibbsMeasure.exists_isProjectiveLimit_finVolMeasure,
      MeasureTheory.GibbsMeasure.eval_localEvents_of_isProjectiveLimit}
    \leanok

    The marginals of Definition \ref{def:cp-finvol} form a projective family. If
    $(E, \mathcal{E})$ is standard Borel, there is a $\mu \in \Prob(\Cfg, \mathcal{F})$ with
    these marginals, and then $\mu(A) = \mu^0(A)$ for every local event $A$.
\end{lemma}
\begin{proof}
    \leanok

    Consistency: $\operatorname{restrict}_{J \subseteq I} \circ \sigma_I = \sigma_J$, so the
    push-forward of the $I$-marginal is the $J$-marginal. Kolmogorov's extension theorem gives
    the projective limit $\mu$; on a cylinder $\{\sigma_\Lambda \in B\}$ it evaluates to the
    $\Lambda$-marginal of $B$, i.e.\ to $\mu^0$.
\end{proof}

\begin{theorem}[Ultrafilter form of Georgii (4.9)]
    \label{thm:cp-ultrafilter-convergence}
    \uses{def:cp-loc-equi, lem:cp-projective}
    \lean{MeasureTheory.GibbsMeasure.exists_tendsto_of_locallyEquicontinuous}
    \leanok

    Let $(E, \mathcal{E})$ be standard Borel and $(\mu_a)_{a \in \iota}$ locally equicontinuous
    along $l$. Then along every ultrafilter $U \leq l$ the family converges locally to some
    $\mu \in \Prob(\Cfg, \mathcal{F})$.
\end{theorem}
\begin{proof}
    \leanok

    Let $\mu$ be the projective limit of Lemma \ref{lem:cp-projective}. On a local event $A$,
    $\mu(A) = \mu^0(A) = \lim_{a \in U}\mu_a(A)$.
\end{proof}

\begin{proposition}[Georgii (4.9)]
    \label{prop:cp-49}
    \uses{def:cp-loc-equi, thm:cp-ultrafilter-convergence, def:cp-standard-borel}
    \lean{MeasureTheory.GibbsMeasure.exists_mapClusterPt_of_locallyEquicontinuous}
    \leanok

    Let $(E, \mathcal{E})$ be standard Borel, with no countability assumption on $S$. Every
    locally equicontinuous family in $\Prob(\Cfg, \mathcal{F})$ indexed along a filter
    $l \neq \bot$, in particular every locally equicontinuous net, has a cluster point in
    $\Prob(\Cfg, \mathcal{F})$.
\end{proposition}
\begin{proof}
    \leanok

    Refine $l$ to an ultrafilter $U$ and apply Theorem \ref{thm:cp-ultrafilter-convergence};
    a limit along an ultrafilter refining $l$ is a cluster point along $l$.
\end{proof}

\subsection*{Uniform domination: Georgii (4.10) and (4.11)}

\begin{definition}[Dominated random fields, Georgii's $\mathcal{K}$ of (4.10)]
    \label{def:cp-dominated}
    \uses{def:cylinder-event}
    \lean{MeasureTheory.GibbsMeasure.dominatedBy}
    \leanok

    Given measures $\nu_\Lambda$ on $(\Cfg, \mathcal{F})$, $\Lambda \in \Fin$, set
    \[\mathcal{K} = \set{\mu \in \Prob(\Cfg, \mathcal{F}) :
      \mu(A) \leq \nu_\Lambda(A) \text{ for all } A \in \mathcal{F}_\Lambda,\ \Lambda \in \Fin}.\]
    Only the restriction of $\nu_\Lambda$ to $\mathcal{F}_\Lambda$ enters.
\end{definition}

\begin{lemma}[Local uniform domination implies local equicontinuity]
    \label{lem:cp-dom-equi}
    \uses{def:cp-loc-equi, def:cp-dominated}
    \lean{MeasureTheory.GibbsMeasure.locallyEquicontinuous_of_eventually_le,
      MeasureTheory.GibbsMeasure.isClosed_dominatedBy}
    \leanok

    Let $\nu_\Lambda$, $\Lambda \in \Fin$, be finite measures on $(\Cfg, \mathcal{F})$. If for
    each $\Lambda \in \Fin$, eventually in $a$, $\mu_a(A) \leq \nu_\Lambda(A)$ for all
    $A \in \mathcal{F}_\Lambda$, then $(\mu_a)_{a \in \iota}$ is locally equicontinuous.
    For arbitrary $\nu_\Lambda$, $\mathcal{K}$ is closed in the topology of local convergence.
\end{lemma}
\begin{proof}
    \leanok

    For $A_m \in \mathcal{F}_\Lambda$ with $A_m \downarrow \emptyset$,
    $\limsup_a \mu_a(A_m) \leq \nu_\Lambda(A_m) \to 0$ by continuity from above of the finite
    measure $\nu_\Lambda$. $\mathcal{K}$ is an intersection of the sets
    $\{\mu : \mu(A) \leq \nu_\Lambda(A)\}$, each closed since $\mu \mapsto \mu(A)$ is
    continuous for a local event $A$.
\end{proof}

\begin{corollary}[Georgii (4.10)]
    \label{cor:cp-410}
    \uses{def:cp-dominated, lem:cp-dom-equi, prop:cp-49, def:cp-standard-borel}
    \lean{MeasureTheory.GibbsMeasure.isCompact_dominatedBy}
    \leanok

    Let $(E, \mathcal{E})$ be standard Borel and $\nu_\Lambda$, $\Lambda \in \Fin$, finite
    measures on $(\Cfg, \mathcal{F})$. Then $\mathcal{K}$ is compact.
\end{corollary}
\begin{proof}
    \leanok

    $\mathcal{K}$ is closed and every family in $\mathcal{K}$ is locally equicontinuous
    (Lemma \ref{lem:cp-dom-equi}), hence converges along any ultrafilter to a limit in
    $\mathcal{K}$ by Theorem \ref{thm:cp-ultrafilter-convergence}. Apply the ultrafilter
    criterion for compactness.
\end{proof}

\begin{example}[Georgii (4.11)(1)]
    \label{ex:cp-411-1}
    \uses{cor:cp-410, def:modif, def:isssd-spec, def:gibbs-meas}
    \lean{MeasureTheory.GibbsMeasure.isCompact_closure_setOf_mem_GP_modification}
    \leanok

    Let $(E, \mathcal{E})$ be standard Borel, $\nu$ a probability measure on $E$, and
    $\gamma = \rho\ISSSD(\nu)$ a modification of the independent specification with
    $\rho_\Lambda \leq C_\Lambda < \infty$ for each $\Lambda \in \Fin$. Then
    $\Gibbs{\rho\ISSSD(\nu)}$ is contained in a set $\mathcal{K}$ as in
    Corollary \ref{cor:cp-410}, and its closure is compact.
\end{example}
\begin{proof}
    \leanok

    For $\mu \in \Gibbs{\rho\ISSSD(\nu)}$ and $A \in \mathcal{F}_\Lambda$,
    $\mu(A) = \int (\rho\ISSSD(\nu))_\Lambda(A \mid \omega)\, \mu(d\omega)
    \leq C_\Lambda\, \nu^S(A)$, since $\ISSSD(\nu)_\Lambda(A \mid \omega) = \nu^S(A)$. So
    $\Gibbs{\rho\ISSSD(\nu)} \subseteq \mathcal{K}$ for $\nu_\Lambda = C_\Lambda \nu^S$, which is
    compact by Corollary \ref{cor:cp-410}.
\end{proof}

\begin{example}[Georgii (4.11)(2)]
    \label{ex:cp-411-2}
    \uses{cor:cp-410}
    \leanok

    Let $E$ be finite with every singleton measurable. Then $\Prob(\Cfg, \mathcal{F})$ is
    compact.
\end{example}
\begin{proof}
    \leanok

    If $\Cfg = \emptyset$ there is nothing to prove. Otherwise push the counting measure on
    $E^\Lambda$ forward along a measurable section $E^\Lambda \to \Cfg$ (extend by a fixed
    configuration off $\Lambda$); a nonempty $B \subseteq E^\Lambda$ has counting measure
    $\geq 1$, so these finite measures $\nu_\Lambda$ dominate every probability measure on
    $\mathcal{F}_\Lambda$, and $\mathcal{K} = \Prob(\Cfg, \mathcal{F})$ is compact by
    Corollary \ref{cor:cp-410}.
\end{proof}

\subsection*{Local equicontinuity of nets of Gibbs distributions}

\begin{theorem}[Georgii (4.12)]
    \label{thm:cp-412}
    \uses{def:cp-loc-equi, def:modif, def:spec, def:isssd-spec}
    \lean{MeasureTheory.GibbsMeasure.locallyEquicontinuous_of_eventually_boundedOn}
    \leanok

    Let $\gamma_0$ be a specification with free measure a probability measure $\mu_0$, i.e.\
    $(\gamma_0)_\Lambda(A \mid \eta) = \mu_0(A)$ for all $\Lambda \in \Fin$,
    $A \in \mathcal{F}_\Lambda$, $\eta \in \Cfg$. Let $\gamma^a = \rho^a\gamma_0$, $a \in \iota$,
    be modifications of $\gamma_0$, $(\nu_a)_{a \in \iota}$ a family in
    $\Prob(\Cfg, \mathcal{F})$, $\Lambda_a \to S$ along $l$ in $\Fin$, and
    $\mu_a = \nu_a \gamma^a_{\Lambda_a}$. Suppose that for each $\Lambda \in \Fin$ and each
    $\varepsilon > 0$ there are $B \in \mathcal{F}$ and $C < \infty$ with
    \[\rho^a_\Lambda \leq C \text{ on } B \text{ eventually in } a,
      \qquad
      \limsup_{a \in l} \mu_a(\Cfg \setminus B) \leq \varepsilon.\]
    Then $(\mu_a)_{a \in \iota}$ is locally equicontinuous, with no assumption on
    $(E, \mathcal{E})$; if $(E, \mathcal{E})$ is standard Borel it has a cluster point by
    Proposition \ref{prop:cp-49}.
\end{theorem}
\begin{proof}
    \leanok

    Let $\Lambda \in \Fin$, $A_m \in \mathcal{F}_\Lambda$ with $A_m \downarrow \emptyset$, and
    choose $B$, $C$ for $\varepsilon/2$. Eventually $\Lambda \subseteq \Lambda_a$, so
    $\mu_a = \mu_a\gamma^a_\Lambda$ by consistency, and
    \[\mu_a(A_m) \leq \mu_a(A_m \cap B) + \mu_a(\Cfg\setminus B)
      \leq C\,\mu_0(A_m) + \mu_a(\Cfg \setminus B),\]
    using $\gamma^a_\Lambda(A_m \cap B \mid \omega) \leq C\,(\gamma_0)_\Lambda(A_m \mid \omega)
    = C\,\mu_0(A_m)$. Hence $\limsup_a \mu_a(A_m) \leq C\,\mu_0(A_m) + \varepsilon/2$, and
    $\mu_0(A_m) \to 0$.
\end{proof}

\begin{lemma}[Georgii (4.12), literal density hypothesis]
    \label{lem:cp-412-limsup}
    \uses{thm:cp-412}
    \lean{MeasureTheory.GibbsMeasure.locallyEquicontinuous_of_limsup_iSup_ne_top}
    \leanok

    Theorem \ref{thm:cp-412} holds with the density hypothesis replaced by
    $\limsup_{a \in l} \sup_{\omega \in B} \rho^a_\Lambda(\omega) < \infty$.
\end{lemma}
\begin{proof}
    \leanok

    Put $C = \limsup_a \sup_{\omega \in B}\rho^a_\Lambda(\omega) + 1 < \infty$. Eventually
    $\sup_B \rho^a_\Lambda < C$; apply Theorem \ref{thm:cp-412}.
\end{proof}

\begin{corollary}[Georgii (4.13), density form]
    \label{cor:cp-413}
    \uses{thm:cp-412}
    \lean{MeasureTheory.GibbsMeasure.locallyEquicontinuous_of_confinement}
    \leanok

    In the setting of Theorem \ref{thm:cp-412}, let $(K_\ell)_{\ell \geq 1}$ be a sequence in
    $\mathcal{E}$ such that
    \begin{enumerate}
        \item $\lim_{\ell \to \infty} \limsup_{a \in l} \mu_a(\sigma_i \notin K_\ell) = 0$ for
              all $i \in S$;
        \item each $\Lambda \in \Fin$ is contained in some $\Delta \in \Fin$ such that, for every
              $\ell$, there is $C < \infty$ with $\rho^a_\Lambda \leq C$ on
              $K_\ell^\Delta \times E^{S \setminus \Delta}$ eventually in $a$.
    \end{enumerate}
    Then $(\mu_a)_{a \in \iota}$ is locally equicontinuous. No positivity or finiteness of
    $\mu_0$-mass of the $K_\ell$ is assumed.
\end{corollary}
\begin{proof}
    \leanok

    Given $\Lambda$ and $\varepsilon > 0$, take $\Delta \supseteq \Lambda$ from (ii). By (i) and
    a union bound over the sites of $\Delta$, choose $\ell$ with
    $\sum_{i \in \Delta}\limsup_a \mu_a(\sigma_i \notin K_\ell) < \varepsilon$. Then
    $B = K_\ell^\Delta \times E^{S\setminus\Delta}$ satisfies
    $\limsup_a \mu_a(\Cfg\setminus B) \leq \varepsilon$, and the densities are eventually
    bounded on $B$ by (ii); apply Theorem \ref{thm:cp-412}.
\end{proof}

\begin{corollary}[Georgii (4.13), Hamiltonian form]
    \label{cor:cp-413-ham}
    \uses{cor:cp-413}
    \lean{Potential.locallyEquicontinuous_of_confinement_hamiltonian}
    \leanok

    Let $S$ be countable, $\nu$ a probability measure on $E$, $\beta \in \R$, and
    $(\Phi^a)_{a \in \iota}$ summable potentials whose Boltzmann factors
    $e^{-\beta H^{\Phi^a}_\Lambda}$ are $\nu$-admissible (Georgii's condition (2.7)). Let
    $\Lambda_a \to S$ along $l$, $\nu_a \in \Prob(\Cfg, \mathcal{F})$, and
    $\mu_a = \nu_a \gamma^{\Phi^a}_{\Lambda_a}$. Let $(K_\ell)_{\ell \geq 1}$ be a sequence in
    $\mathcal{E}$ with $\nu(K_\ell) > 0$ for all $\ell$, satisfying (i) of
    Corollary \ref{cor:cp-413}, and such that each $\Lambda \in \Fin$ is contained in some
    $\Delta \in \Fin$ with, for every $\ell$, a $c \in \R$ such that
    \[|H^{\Phi^a}_\Lambda| \leq c \text{ on } K_\ell^\Delta \times E^{S\setminus\Delta}
      \text{ eventually in } a.\]
    Then $(\mu_a)_{a \in \iota}$ is locally equicontinuous.
\end{corollary}
\begin{proof}
    \leanok

    From $|H^{\Phi^a}_\Lambda| \leq c$ on the box, the Boltzmann factor is
    $\leq e^{|\beta|c}$ there and the partition function is
    $\geq e^{-|\beta|c}\nu(K_\ell)^{|\Lambda|} > 0$, so
    $\rho^{\Phi^a}_\Lambda \leq e^{|\beta|c}\big/\big(e^{-|\beta|c}\nu(K_\ell)^{|\Lambda|}\big)$
    on the box. This is (ii) of Corollary \ref{cor:cp-413}.
\end{proof}

\begin{remark}[Georgii (4.14)]
    \label{rmk:cp-414}
    \uses{cor:cp-413-ham}

    (1) The hypotheses of Corollary \ref{cor:cp-413-ham} hold when
    $\limsup_{a} \|\Phi^a\|_i < \infty$ for all $i \in S$: take $K_\ell = E$ and
    $\Delta = \Lambda$.
    (2) The Hamiltonian bound also follows from the existence of a finite-range potential
    $\Phi$ with $\limsup_{a} \|\Phi^a - \Phi\|_i < \infty$ for all $i$ and each $\Phi_A$
    bounded on $K_\ell^S$.
\end{remark}

\subsection*{Georgii (4.15): from cluster point to convergent subsequence}

\begin{definition}[Georgii's $\nu$, $h_{n,\Lambda}$, $\mathcal{A}_\Lambda$, $\mathcal{A}^0$]
    \label{def:cp-dens}
    \lean{MeasureTheory.GibbsMeasure.domMeasure,
      MeasureTheory.GibbsMeasure.localDensity,
      MeasureTheory.GibbsMeasure.localDens,
      MeasureTheory.GibbsMeasure.densAlgebra,
      MeasureTheory.GibbsMeasure.densSeedsLocal,
      MeasureTheory.GibbsMeasure.densSeeds}
    \leanok

    For a sequence $(\mu_n)_{n \geq 1}$ and a candidate limit $\mu$, put
    $\nu = \mu + \sum_{n \geq 1} 2^{-n}\mu_n$; let $h_{n,\Lambda}$ be a Radon--Nikodym density
    of $\mu_n|\mathcal{F}_\Lambda$ with respect to $\nu|\mathcal{F}_\Lambda$, let
    $\mathcal{A}_\Lambda \subseteq \mathcal{F}_\Lambda$ be the $\sigma$-algebra generated by the
    $h_{n,\Lambda}$, and let $\mathcal{A}^0_\Lambda$ be the family of finite intersections of
    preimages, under the $h_{n,\Lambda}$, of a fixed countable generating family of the Borel
    sets of $[0,\infty]$. Put $\mathcal{A}^0 = \bigcup_{\Lambda \in \Fin} \mathcal{A}^0_\Lambda$.
\end{definition}

\begin{lemma}[$\mathcal{A}^0$ is a countable $\pi$-system of local events generating the
  $\mathcal{A}_\Lambda$]
    \label{lem:cp-dens-countable}
    \uses{def:cp-dens}
    \lean{MeasureTheory.GibbsMeasure.countable_densSeedsLocal,
      MeasureTheory.GibbsMeasure.isPiSystem_densSeedsLocal,
      MeasureTheory.GibbsMeasure.generateFrom_densSeedsLocal,
      MeasureTheory.GibbsMeasure.countable_densSeeds,
      MeasureTheory.GibbsMeasure.densSeeds_subset_localEvents}
    \leanok

    Each $\mathcal{A}^0_\Lambda$ is a countable $\pi$-system generating $\mathcal{A}_\Lambda$,
    and, for $S$ countable, $\mathcal{A}^0$ is a countable family of local events.
\end{lemma}
\begin{proof}
    \leanok

    Finite intersections of a countable family form a countable $\pi$-system generating the same
    $\sigma$-algebra; the preimages of a countable generating family of the Borel sets of
    $[0,\infty]$ under the countably many $h_{n,\Lambda}$ generate $\mathcal{A}_\Lambda$. Each
    generator is $\mathcal{F}_\Lambda$-measurable, hence a local event. For $S$ countable the
    union over $\Lambda \in \Fin$ is countable.
\end{proof}

\begin{lemma}[Transfer from $\mathcal{A}_\Lambda$ back to $\mathcal{F}_\Lambda$]
    \label{lem:cp-transfer}
    \uses{def:cp-dens, def:lebesgue-cond-exp}
    \lean{MeasureTheory.GibbsMeasure.condDensity,
      MeasureTheory.GibbsMeasure.setLIntegral_eq_lintegral_mul_condDensity,
      MeasureTheory.GibbsMeasure.measure_eq_lintegral_condDensity,
      MeasureTheory.GibbsMeasure.tendsto_lintegral_of_tendsto_setwise}
    \leanok

    Let $A \in \mathcal{F}_\Lambda$ and let $f_A$ be the version of
    $\nu(A \mid \mathcal{A}_\Lambda)$ given by a Radon--Nikodym density of
    $\nu(A \cap \cdot)|\mathcal{A}_\Lambda$ with respect to $\nu|\mathcal{A}_\Lambda$, truncated
    at $1$. Then $\rho(A) = \rho(f_A)$ for every measure $\rho \ll \nu$ whose density on
    $\mathcal{F}_\Lambda$ is $\mathcal{A}_\Lambda$-measurable, in particular for each $\mu_n$
    and for $\mu$. Setwise convergence on $\mathcal{A}_\Lambda$ of probability measures
    implies convergence of the integrals of every $\mathcal{A}_\Lambda$-measurable function
    bounded by $1$.
\end{lemma}
\begin{proof}
    \leanok

    With $h = h_{n,\Lambda}$, $\mu_n(A) = \nu(h 1_A) = \nu(h f_A) = \mu_n(f_A)$, the middle
    equality by the defining property of $f_A$ and $\mathcal{A}_\Lambda$-measurability of $h$.
    For the last assertion approximate $f$ from below by the layer sums
    $N^{-1}\sum_{j} 1_{\{f > j/N\}}$, whose integrals converge by hypothesis, with uniform error
    $N^{-1}$.
\end{proof}

\begin{lemma}[Dynkin step of Georgii (4.15)]
    \label{lem:cp-dynkin}
    \uses{def:cp-loc-equi, def:cp-dens, lem:cp-dens-countable, lem:cp-transfer}
    \lean{MeasureTheory.GibbsMeasure.tendsto_of_tendsto_densSeeds}
    \leanok

    Let $(\mu_n)_{n\geq 1}$ be locally equicontinuous and $\mu$ a probability measure. If a
    subsequence $(\mu_{n_k})$ satisfies $\mu_{n_k}(A) \to \mu(A)$ for every $A \in \mathcal{A}^0$,
    then $\mu_{n_k}(A) \to \mu(A)$ for every local event $A$.
\end{lemma}
\begin{proof}
    \leanok

    Fix $\Lambda$ and let $\mathcal{A}^1_\Lambda$ be the sets of $\mathcal{A}_\Lambda$ on which
    convergence holds. It contains the $\pi$-system $\mathcal{A}^0_\Lambda$ and is a Dynkin
    system: for a countable disjoint union $B = \bigcup_\ell B_\ell$, local equicontinuity on
    the tails $B \setminus C_m \downarrow \emptyset$, $C_m = \bigcup_{\ell \leq m}B_\ell$, gives
    $\limsup_k \mu_{n_k}(B) \leq \mu(B)+\varepsilon$ and
    $\liminf_k \mu_{n_k}(B) \geq \mu(B)-\varepsilon$. Hence
    $\mathcal{A}^1_\Lambda = \mathcal{A}_\Lambda$, and Lemma \ref{lem:cp-transfer} carries the
    convergence to $\mathcal{F}_\Lambda$ via $\mu_n(A) = \mu_n(f_A)$.
\end{proof}

\begin{proposition}[Georgii (4.15)]
    \label{prop:cp-415}
    \uses{def:cp-loc-equi, lem:cp-dynkin, lem:cp-dens-countable}
    \lean{MeasureTheory.GibbsMeasure.exists_subseq_tendsto_of_mapClusterPt,
      MeasureTheory.GibbsMeasure.exists_strictMono_tendsto_of_mapClusterPt}
    \leanok

    Let $S$ be countable, with no assumption on $(E, \mathcal{E})$. If
    $\mu \in \Prob(\Cfg, \mathcal{F})$ is a cluster point of a locally equicontinuous sequence
    $(\mu_n)_{n \geq 1}$ in $\Prob(\Cfg, \mathcal{F})$, then some subsequence
    $(\mu_{n_k})_{k\geq 1}$ converges to $\mu$.
\end{proposition}
\begin{proof}
    \leanok

    Enumerate $\mathcal{A}^0 = \{A_1, A_2, \dots\}$ (Lemma \ref{lem:cp-dens-countable}). The map
    $\rho \mapsto (\rho(A_\ell))_{\ell \geq 1}$ into the metrizable space
    $[0,\infty]^{\mathbb{N}}$ is continuous for local convergence, so there is a strictly
    increasing $(n_k)$ with $\mu_{n_k}(A_\ell) \to \mu(A_\ell)$ for every $\ell$.
    Lemma \ref{lem:cp-dynkin} extends this to all local events.
\end{proof}

\begin{corollary}[Georgii (4.9) and (4.15) combined]
    \label{cor:cp-49-415}
    \uses{prop:cp-49, prop:cp-415}
    \lean{MeasureTheory.GibbsMeasure.exists_subseq_tendsto_of_locallyEquicontinuous}
    \leanok

    Let $(E, \mathcal{E})$ be standard Borel and $S$ countable. Every locally equicontinuous
    sequence in $\Prob(\Cfg, \mathcal{F})$ has a subsequence converging in the topology of
    local convergence.
\end{corollary}
\begin{proof}
    \leanok

    Proposition \ref{prop:cp-49} gives a cluster point; Proposition \ref{prop:cp-415} makes it
    the limit of a subsequence.
\end{proof}
